% sec_precision_potential.tex -- Poisson smoothing, entropy, and reconstruction
\section{Poisson Smoothing, Entropy, and Reconstruction}%
\label{sec:precision-potential}

\subsection{Entropy in Cayley coordinates}

The logarithmic Jacobian of the Cayley coordinate is
\begin{equation}\label{eq:precision-potential}
\mathcal P(x):=-\log J(x)=\log\pi+\log(1+x^{2}).
\end{equation}

\begin{theorem}[Entropy identity in Cayley coordinates]%
\label{thm:precision-ledger}
Let $X$ be a real-valued random variable with density $f_X$ such that
$h_{\mathrm{diff}}(X)>-\infty$ and $\EE[\mathcal P(X)]<\infty$.  Set
$U:=\upsilon(X)\in\UU$.  Then
\begin{equation}\label{eq:entropy-ledger}
h_{\mathrm{diff}}(X)
=h_{\mathrm{diff}}(U)+\EE[\mathcal P(X)].
\end{equation}
\end{theorem}

\begin{proof}
The density of $U$ is
\[
f_U(u)=f_X(\tan(\pi u))\,\pi(1+\tan^2(\pi u)).
\]
The standard change-of-variables formula for differential entropy gives
\[
h_{\mathrm{diff}}(X)
=h_{\mathrm{diff}}(U)+\EE\!\left[\log\frac{1}{\abs{\upsilon'(X)}}\right]
=h_{\mathrm{diff}}(U)+\EE[\mathcal P(X)].
\]
\end{proof}

\subsection{Poisson smoothing}

For $t>0$ and $x\in\RR$ define
\begin{equation}\label{eq:poisson-kernel}
P_t(x):=\frac{1}{\pi}\frac{t}{x^{2}+t^{2}}.
\end{equation}

\begin{proposition}[Basic identities for the Poisson kernels]%
\label{prop:poisson-identities}
For $s,t>0$ one has
\begin{equation}\label{eq:poisson-basic}
P_t(x)=\frac{1}{t}P_1\!\left(\frac{x}{t}\right),
\qquad
\widehat{P_t}(\xi)=e^{-t\abs{\xi}},
\qquad
P_t*P_s=P_{t+s}.
\end{equation}
In particular,
\[
\mu=P_1(x)\dd x,
\qquad
P_t*\mu=P_{t+1}.
\]
\end{proposition}

\begin{proof}
The scaling identity is immediate from the definition.  The Fourier
transform of~$P_1$ is $e^{-\abs{\xi}}$
\cite[Chapter~III]{Stein1971}; scaling yields
$\widehat{P_t}(\xi)=e^{-t\abs{\xi}}$.  The convolution law follows by
multiplying Fourier transforms.
\end{proof}

\begin{proposition}[Standard facts on the Poisson semigroup]%
\label{thm:poisson-semigroup-law}
For $1\le p\le\infty$ and $t>0$ define
$T_t f:=P_t*f$.  Then:
\begin{enumerate}[label=(\roman*)]
\item $\|T_t f\|_{L^p(\RR)}\le \|f\|_{L^p(\RR)}$ for every $f\in L^p$;
\item $T_{t+s}=T_tT_s$ for all $s,t>0$;
\item for $1\le p<\infty$, $T_t f\to f$ in $L^p(\RR)$ as $t\downarrow 0$;
\item on $L^2(\RR)$, $\widehat{T_t f}(\xi)=e^{-t\abs{\xi}}\widehat f(\xi)$
      and the generator is $-|D|$.
\end{enumerate}
\end{proposition}

\begin{proof}
Since $\int_{\RR}P_t(x)\dd x=1$ and $P_t\ge 0$, Young's inequality gives
(i).  The semigroup law is Proposition~\ref{prop:poisson-identities}.
Strong continuity for $1\le p<\infty$ follows because $\{P_t\}_{t>0}$ is
an approximate identity.  The $L^2$-multiplier statement is exactly
\eqref{eq:poisson-basic}; the generator follows from
\[
\frac{e^{-t\abs{\xi}}-1}{t}\longrightarrow -\abs{\xi}
\qquad(t\downarrow 0)
\]
in the Fourier representation.
\end{proof}

\begin{remark}
The preceding proposition is standard and is included only to fix the
normalization of the Poisson kernel and of the generator $-|D|$ used in
the sequel.
\end{remark}

\subsection{Large-time comparison with a translated Poisson profile}

\begin{proposition}[Second-order approximation]%
\label{thm:poisson-second-order}
Let $\nu$ be a probability measure on~$\RR$ with mean
\[
\bar\gamma:=\int_{\RR}\gamma\dd\nu(\gamma),
\]
variance
\[
\sigma^{2}:=\int_{\RR}(\gamma-\bar\gamma)^{2}\dd\nu(\gamma),
\]
and finite centred fourth moment.  Define
\[
h_t:=P_t*\nu,
\qquad
g_t(x):=P_t(x-\bar\gamma).
\]
Let
\[
U(y):=\frac{3y^{2}-1}{(1+y^{2})^{2}}.
\]
Then there exists a constant $C_\nu>0$ such that for every $t\ge 1$,
\begin{equation}\label{eq:poisson-profile}
\sup_{x\in\RR}\left|
\frac{h_t(x)}{g_t(x)}
-1
-\frac{\sigma^{2}}{t^{2}}U\!\left(\frac{x-\bar\gamma}{t}\right)
\right|
\le \frac{C_\nu}{t^{3}}.
\end{equation}
Consequently,
\begin{equation}\label{eq:poisson-l1-rate}
\|h_t-g_t\|_{L^{1}(\RR)}=O(t^{-2}),
\qquad
\|h_t-g_t\|_{L^\infty(\RR)}=O(t^{-3}).
\end{equation}
Moreover,
\begin{equation}\label{eq:poisson-kl-rate}
D_{\mathrm{KL}}(h_t\|g_t)=O(t^{-4}).
\end{equation}
\end{proposition}

\begin{proof}
Write $\Delta:=\gamma-\bar\gamma$ and $x=\bar\gamma+ty$.  Then
\[
\frac{P_t(x-\gamma)}{g_t(x)}
=\frac{1+y^{2}}{1+(y-\Delta/t)^{2}}.
\]
Set $\varepsilon:=\Delta/t$.  A direct algebraic computation gives the
exact identity
\begin{align*}
\frac{1+y^{2}}{1+(y-\varepsilon)^{2}}
&=1+\frac{2y}{1+y^{2}}\varepsilon
+\frac{3y^{2}-1}{(1+y^{2})^{2}}\varepsilon^{2} \\
&\quad
+\frac{4y(y^{2}-1)\varepsilon^{3}-(3y^{2}-1)\varepsilon^{4}}
{(1+y^{2})^{2}(1+(y-\varepsilon)^{2})}.
\end{align*}
Since
\[
\frac{4\abs{y}\abs{y^{2}-1}}{(1+y^{2})^{2}}\le 2,
\qquad
\frac{\abs{3y^{2}-1}}{(1+y^{2})^{2}}\le 1,
\]
the remainder is bounded by $2\abs{\varepsilon}^{3}+\abs{\varepsilon}^{4}$
uniformly in~$y$.  Taking expectation and using $\EE[\Delta]=0$ yields
\[
\frac{h_t(\bar\gamma+ty)}{g_t(\bar\gamma+ty)}
=1+\frac{\sigma^{2}}{t^{2}}U(y)+R_t(y),
\]
with
\[
\sup_{y\in\RR}\abs{R_t(y)}
\le \frac{2\EE[\abs{\Delta}^{3}]+\EE[\abs{\Delta}^{4}]}{t^{3}}
\]
for $t\ge 1$.  This proves \eqref{eq:poisson-profile}.

The function~$U$ is bounded, so
$\sup_x\abs{h_t(x)/g_t(x)-1}=O(t^{-2})$.  Since $\int_{\RR}g_t=1$,
\[
\|h_t-g_t\|_{L^1}
\le \sup_{x}\abs{\frac{h_t(x)}{g_t(x)}-1}\int_{\RR}g_t(x)\dd x
=O(t^{-2}).
\]
Also $\|g_t\|_{L^\infty}=1/(\pi t)$, so
\[
\|h_t-g_t\|_{L^\infty}
\le \|g_t\|_{L^\infty}\sup_{x}\abs{\frac{h_t(x)}{g_t(x)}-1}
=O(t^{-3}).
\]
Let
\[
\delta_t(x):=\frac{h_t(x)}{g_t(x)}-1.
\]
The bound \eqref{eq:poisson-profile} implies that
$\|\delta_t\|_{L^\infty}=O(t^{-2})$, so for large $t$ one has
$\|\delta_t\|_{L^\infty}\le 1/2$.  Since
$\int g_t\delta_t\dd x=\int(h_t-g_t)\dd x=0$ and
\[
(1+s)\log(1+s)-s\le 2s^{2}
\qquad(\abs{s}\le 1/2),
\]
it follows that
\[
D_{\mathrm{KL}}(h_t\|g_t)
=\int_{\RR}g_t(x)\bigl((1+\delta_t(x))\log(1+\delta_t(x))-\delta_t(x)\bigr)\dd x
\le 2\|\delta_t\|_{L^\infty}^{2}.
\]
Hence \eqref{eq:poisson-kl-rate} holds.
\end{proof}

\subsection{Cayley--Chebyshev mode calculus}

For $\varepsilon\in\RR$ and $y\in\RR$ set
\begin{equation}\label{eq:cayley-comparison-kernel}
R_\varepsilon(y):=\frac{1+y^2}{1+(y-\varepsilon)^2}.
\end{equation}
For fixed $y$, the Taylor expansion of $R_\varepsilon(y)$ at
$\varepsilon=0$ defines coefficients $u_n(y)$ through
\begin{equation}\label{eq:cayley-mode-series}
R_\varepsilon(y)=\sum_{n=0}^{\infty}u_n(y)\varepsilon^n,
\qquad
|\varepsilon|<\sqrt{1+y^2}.
\end{equation}

\begin{theorem}[Cayley--Chebyshev mode formula]%
\label{thm:cayley-chebyshev-mode}
Let $U_n$ denote the Chebyshev polynomials of the second kind.  If
$y=\tan\theta$ with $\theta\in(-\pi/2,\pi/2)$, then for every $n\ge 0$,
\begin{equation}\label{eq:cayley-chebyshev-mode}
u_n(\tan\theta)=\cos^n\theta\,U_n(\sin\theta).
\end{equation}
For $n\ge 1$ this is equivalent to
\begin{equation}\label{eq:cayley-chebyshev-trig}
u_n(\tan\theta)
=\cos^{n-1}\theta\,
\sin\!\Bigl((n+1)\Bigl(\frac{\pi}{2}-\theta\Bigr)\Bigr).
\end{equation}
Consequently
\begin{equation}\label{eq:cayley-mode-parity}
u_n(-y)=(-1)^n u_n(y)
\qquad(n\ge 0).
\end{equation}
Moreover, for every integer $m\ge 1$,
\begin{equation}\label{eq:cayley-mode-even-fourier}
u_{2m}(\tan\theta)
=(-1)^m 2^{-(2m-1)}
\sum_{j=0}^{2m-1}\binom{2m-1}{j}\cos(2(j+1)\theta),
\end{equation}
and for every integer $m\ge 0$,
\begin{equation}\label{eq:cayley-mode-odd-fourier}
u_{2m+1}(\tan\theta)
=(-1)^m 2^{-2m}
\sum_{j=0}^{2m}\binom{2m}{j}\sin(2(j+1)\theta).
\end{equation}
\end{theorem}

\begin{proof}
With $y=\tan\theta$ one has
\[
R_\varepsilon(\tan\theta)
=\frac{1}{1-2(\varepsilon\cos\theta)\sin\theta+(\varepsilon\cos\theta)^2}.
\]
The generating function
\[
\sum_{n=0}^{\infty}U_n(x)r^n=\frac{1}{1-2xr+r^2}
\]
therefore yields \eqref{eq:cayley-chebyshev-mode} with
$x=\sin\theta$ and $r=\varepsilon\cos\theta$.  Since
\[
U_n(\sin\theta)
=\frac{\sin((n+1)(\pi/2-\theta))}{\cos\theta},
\]
formula \eqref{eq:cayley-chebyshev-trig} follows for $n\ge 1$.  The
parity relation \eqref{eq:cayley-mode-parity} is immediate from
\eqref{eq:cayley-chebyshev-trig}.

For the finite Fourier expansions, write
\[
u_{2m}(\tan\theta)=(-1)^m\cos^{2m-1}\theta\,\cos((2m+1)\theta),
\]
\[
u_{2m+1}(\tan\theta)=(-1)^m\cos^{2m}\theta\,\sin((2m+2)\theta),
\]
and use
\[
1+e^{2\ii\theta}=2e^{\ii\theta}\cos\theta.
\]
Taking real and imaginary parts after the binomial expansion gives
\eqref{eq:cayley-mode-even-fourier} and
\eqref{eq:cayley-mode-odd-fourier}.
\end{proof}

\begin{proposition}[Finite mode truncations]%
\label{prop:cayley-mode-remainder}
For every integer $N\ge 0$ there exists a rational function
$v_N:\RR^2\to\RR$, bounded on~$\RR^2$, such that
\begin{equation}\label{eq:cayley-mode-remainder}
R_\varepsilon(y)
=\sum_{n=0}^{N}u_n(y)\varepsilon^n+\varepsilon^{N+1}v_N(y,\varepsilon)
\qquad(y,\varepsilon\in\RR).
\end{equation}
\end{proposition}

\begin{proof}
For every $x$ and $r$ one has the finite-sum identity
\[
\sum_{n=0}^{N}U_n(x)r^n
=\frac{1-r^{N+1}U_{N+1}(x)+r^{N+2}U_N(x)}{1-2xr+r^2}.
\]
Substituting $x=\sin\theta$ and $r=\varepsilon\cos\theta$ gives
\eqref{eq:cayley-mode-remainder} with
\[
v_N(\tan\theta,\varepsilon)
=\frac{\cos^{N+1}\theta\,U_{N+1}(\sin\theta)
-\varepsilon\cos^{N+2}\theta\,U_N(\sin\theta)}
{1-2(\varepsilon\cos\theta)\sin\theta+(\varepsilon\cos\theta)^2}.
\]
This expression is rational in $(y,\varepsilon)$.  Since the denominator
is positive on $\RR^2$ and the numerator has lower degree than the
denominator, $v_N$ is bounded on~$\RR^2$.
\end{proof}

For later use we record the first modes explicitly:
\begin{equation}\label{eq:cayley-first-modes}
u_2(y)=\frac{3y^2-1}{(1+y^2)^2},
\qquad
u_3(y)=\frac{4y(y^2-1)}{(1+y^2)^3},
\qquad
u_4(y)=\frac{5y^4-10y^2+1}{(1+y^2)^4},
\end{equation}
\begin{equation}\label{eq:cayley-next-modes}
u_5(y)=\frac{6y^5-20y^3+6y}{(1+y^2)^5},
\qquad
u_6(y)=\frac{7y^6-35y^4+21y^2-1}{(1+y^2)^6}.
\end{equation}

\subsection{Entropy asymptotics and rigidity}

\begin{theorem}[Odd orders vanish in the entropy expansion]%
\label{thm:poisson-kl-even-orders}
Let $N\ge 2$ and let $\nu$ be a probability measure on~$\RR$ with mean
$\bar\gamma$ and finite centred $(2N+1)$st absolute moment.  Set
\[
\Delta:=\gamma-\bar\gamma,
\qquad
\mu_n:=\int_{\RR}\Delta^n\dd\nu(\gamma),
\]
and keep the notation
\[
h_t:=P_t*\nu,
\qquad
g_t(x):=P_t(x-\bar\gamma).
\]
Then there exist real coefficients $c_4,c_6,\dots,c_{2N}$ such that
\begin{equation}\label{eq:poisson-kl-even-orders}
D_{\mathrm{KL}}(h_t\|g_t)
=\sum_{m=2}^{N}c_{2m}\,t^{-2m}+o(t^{-2N})
\qquad(t\to\infty).
\end{equation}
In particular, every odd power $t^{-(2m+1)}$ has zero coefficient.
\end{theorem}

\begin{proof}
With $x=\bar\gamma+ty$ one has
\[
\delta_t(x):=\frac{h_t(x)}{g_t(x)}-1
=\int_{\RR}\bigl(R_{\Delta/t}(y)-1\bigr)\dd\nu(\gamma).
\]
Applying Proposition~\ref{prop:cayley-mode-remainder} with $N=2N$ gives
\[
\delta_t(\bar\gamma+ty)
=\sum_{n=2}^{2N}\frac{\mu_n}{t^n}u_n(y)+O(t^{-(2N+1)})
\]
uniformly in $y\in\RR$.  Since $\|\delta_t\|_{L^\infty}=O(t^{-2})$ and
\[
\Phi(z):=(1+z)\log(1+z)-z
=\sum_{k=2}^{\infty}\frac{(-1)^k}{k(k-1)}z^k
\qquad(|z|<1),
\]
one may truncate the power series at order~$N$ and write
\[
\Phi(\delta_t)
=\sum_{k=2}^{N}\frac{(-1)^k}{k(k-1)}\delta_t^k+O(t^{-2N-2})
\]
uniformly in $y$.  Expanding each finite power $\delta_t^k$ with the
uniform remainder above yields an asymptotic expansion of the form
\[
D_{\mathrm{KL}}(h_t\|g_t)
=\int_{\RR}\frac{\Phi(\delta_t(\bar\gamma+ty))}{\pi(1+y^2)}\dd y
\]
\[
=\sum_{m=2}^{N} c_{2m}'\,t^{-m}+o(t^{-2N}),
\]
where each coefficient $c_m'$ is a finite linear combination of integrals
of the form
\[
\int_{\RR}\frac{u_{n_1}(y)\cdots u_{n_k}(y)}{\pi(1+y^2)}\dd y,
\qquad
n_1+\cdots+n_k=m.
\]
The weight $(\pi(1+y^2))^{-1}$ is even, and
\eqref{eq:cayley-mode-parity} implies that the integrand above is odd
unless $n_1+\cdots+n_k$ is even.  Hence $c_m'=0$ for odd~$m$.  Writing
$c_{2m}:=c_{2m}'$ gives \eqref{eq:poisson-kl-even-orders}.
\end{proof}

\begin{theorem}[Eighth-order entropy expansion]%
\label{thm:poisson-kl-eighth}
Let $\nu$ be a probability measure on~$\RR$ with mean~$\bar\gamma$,
variance $\sigma^2>0$, and finite centred seventh absolute moment.  Then,
as $t\to\infty$,
\begin{equation}\label{eq:poisson-kl-eighth}
\begin{split}
D_{\mathrm{KL}}(h_t\|g_t)
=\frac{\sigma^4}{8t^4}
+\frac{\sigma^6+6\mu_3^2-8\sigma^2\mu_4}{64\,t^6}
\\
\quad
+\frac{3\sigma^8-12\sigma^4\mu_4+9\sigma^2\mu_3^2}{256\,t^8}
+\frac{12\sigma^2\mu_6-30\mu_3\mu_5+20\mu_4^2}{256\,t^8}
+o(t^{-8}).
\end{split}
\end{equation}
\end{theorem}

\begin{proof}
With $x=\bar\gamma+ty$, Proposition~\ref{prop:cayley-mode-remainder} with
$N=6$ gives
\[
\delta_t(\bar\gamma+ty)
=\frac{\sigma^2}{t^2}u_2(y)
+\frac{\mu_3}{t^3}u_3(y)
+\frac{\mu_4}{t^4}u_4(y)
+\frac{\mu_5}{t^5}u_5(y)
+\frac{\mu_6}{t^6}u_6(y)
+O(t^{-7})
\]
uniformly in $y\in\RR$.  Since $\delta_t=O(t^{-2})$, one has
\[
\Phi(\delta_t)
=\frac12\delta_t^2-\frac16\delta_t^3+\frac1{12}\delta_t^4+O(t^{-10}).
\]
Also $g_t(\bar\gamma+ty)\dd(\bar\gamma+ty)=\pi^{-1}(1+y^2)^{-1}\dd y$.
Write
\[
a_2:=\sigma^2u_2,\qquad a_3:=\mu_3u_3,\qquad a_4:=\mu_4u_4,
\qquad a_5:=\mu_5u_5,\qquad a_6:=\mu_6u_6,
\]
so that $\delta_t=t^{-2}a_2+t^{-3}a_3+t^{-4}a_4+t^{-5}a_5+t^{-6}a_6+O(t^{-7})$.
Up to order $t^{-8}$ one therefore has
\[
\frac12\delta_t^2
=\frac{a_2^2}{2t^4}
+\frac{a_2a_4+\frac12 a_3^2}{t^6}
+\frac{a_2a_6+a_3a_5+\frac12 a_4^2}{t^8}
+O(t^{-9}),
\]
\[
-\frac16\delta_t^3
=-\frac{a_2^3}{6t^6}
-\frac{\frac12 a_2^2a_4+\frac12 a_2a_3^2}{t^8}
+O(t^{-9}),
\]
\[
\frac1{12}\delta_t^4
=\frac{a_2^4}{12t^8}+O(t^{-9}).
\]
After integration against $\pi^{-1}(1+y^2)^{-1}\dd y$, the weighted
identities from Appendix~\ref{app:entropy-integrals} give
\[
\frac12\int_{\RR}\frac{u_2(y)^2}{\pi(1+y^2)}\dd y=\frac18,
\]
\[
\int_{\RR}\frac{\sigma^2\mu_4u_2u_4+\frac12\mu_3^2u_3^2-\frac16\sigma^6u_2^3}
{\pi(1+y^2)}\dd y
=\frac{\sigma^6+6\mu_3^2-8\sigma^2\mu_4}{64},
\]
and
\begin{align*}
\int_{\RR}
\frac{\sigma^2\mu_6u_2u_6+\mu_3\mu_5u_3u_5+\frac12\mu_4^2u_4^2}
{\pi(1+y^2)}\dd y
&=\frac{12\sigma^2\mu_6-30\mu_3\mu_5+20\mu_4^2}{256},
\\
\int_{\RR}
\frac{-\frac12\sigma^4\mu_4u_2^2u_4-\frac12\sigma^2\mu_3^2u_2u_3^2
+\frac1{12}\sigma^8u_2^4}{\pi(1+y^2)}\dd y
&=\frac{3\sigma^8-12\sigma^4\mu_4+9\sigma^2\mu_3^2}{256}.
\end{align*}
This proves \eqref{eq:poisson-kl-eighth}.  The coefficient of order
$t^{-6}$ is the first non-trivial even coefficient furnished by
Theorem~\ref{thm:poisson-kl-even-orders}; the order-$t^{-8}$ term is the
next level in the same expansion.
\end{proof}

\begin{theorem}[A two-level defect ladder]%
\label{thm:poisson-defect-ladder}
Under the hypotheses of Theorem~\ref{thm:poisson-kl-eighth}, define
\begin{align*}
C_6&:=\frac{\sigma^6+6\mu_3^2-8\sigma^2\mu_4}{64}, \\
C_8&:=
\frac{3\sigma^8-12\sigma^4\mu_4+9\sigma^2\mu_3^2+12\sigma^2\mu_6}{256}
+\frac{-30\mu_3\mu_5+20\mu_4^2}{256}.
\end{align*}
Let $X:=(\gamma-\bar\gamma)/\sigma$ on the probability space
$(\RR,\nu)$, set
\[
\beta_3:=\EE[X^3],
\qquad
Q_2:=X^2-1-\beta_3X,
\]
and define
\[
Q_3:=XQ_2-\|Q_2\|_{L^2(\nu)}^2X-\frac{\beta_3}{4}Q_2.
\]
Then
\begin{equation}\label{eq:defect-c6}
\frac{64}{\sigma^6}C_6
=-7-2\beta_3^2-8\|Q_2\|_{L^2(\nu)}^2,
\end{equation}
and
\begin{equation}\label{eq:defect-c8}
\frac{256}{\sigma^8}C_8
=23
+12\|Q_3\|_{L^2(\nu)}^2
+32\|Q_2\|_{L^2(\nu)}^4
+\frac{61}{4}\beta_3^2\|Q_2\|_{L^2(\nu)}^2
+52\|Q_2\|_{L^2(\nu)}^2
+2\beta_3^4
+13\beta_3^2.
\end{equation}
In particular
\begin{equation}\label{eq:defect-c8-lower}
C_8\ge \frac{23}{256}\sigma^8,
\end{equation}
and equality holds if and only if
\[
\nu=\frac12\delta_{\bar\gamma+\sigma}
+\frac12\delta_{\bar\gamma-\sigma}.
\]
\end{theorem}

\begin{proof}
Write
\[
\kappa:=\|Q_2\|_{L^2(\nu)}^2=\EE[Q_2^2],
\qquad
L:=\EE[XQ_2^2],
\qquad
M:=\EE[Q_2^3].
\]
Since $Q_2=X^2-1-\beta_3X$, one has
\[
\kappa=\beta_4-1-\beta_3^2,
\]
and therefore
\[
\frac{64}{\sigma^6}C_6
=1+6\beta_3^2-8\beta_4
=-7-2\beta_3^2-8\kappa,
\]
which is exactly \eqref{eq:defect-c6}.

The identity $X^2=1+\beta_3X+Q_2$ yields
\[
\beta_5=\beta_3(\beta_3^2+2+2\kappa)+L,
\]
\[
\beta_6=1+3\beta_3^2+3\kappa+\beta_3^4+3\beta_3^2\kappa+3\beta_3L+M.
\]
Hence, if
\[
P_8:=3-12\beta_4+9\beta_3^2+12\beta_6-30\beta_3\beta_5+20\beta_4^2,
\]
then a direct substitution gives
\[
P_8
=23+20\kappa^2+16\beta_3^2\kappa+64\kappa+2\beta_3^4+13\beta_3^2+6\beta_3L+12M.
\]
On the other hand,
\[
\|Q_3\|_{L^2(\nu)}^2
=M+\frac{\beta_3}{2}L+\kappa-\kappa^2+\frac{\beta_3^2}{16}\kappa.
\]
Substituting this identity into the preceding expression for $P_8$
yields \eqref{eq:defect-c8}.  Since
$P_8=256C_8/\sigma^8$, inequality \eqref{eq:defect-c8-lower} follows.

If equality holds in \eqref{eq:defect-c8-lower}, then every non-negative
term on the right-hand side of \eqref{eq:defect-c8} vanishes.  In
particular $\beta_3=0$ and $\|Q_2\|_{L^2(\nu)}=0$, so $X^2=1$
$\nu$-almost everywhere.  Since $\EE[X]=0$ and $\EE[X^2]=1$, it follows
that $X=\pm 1$ with equal weights.  Thus
\[
\nu=\frac12\delta_{\bar\gamma+\sigma}
+\frac12\delta_{\bar\gamma-\sigma}.
\]
The converse is immediate.
\end{proof}

\begin{theorem}[Defect amplification]\label{thm:poisson-defect-amplification}
Under the hypotheses of Theorem~\ref{thm:poisson-kl-eighth}, define
\[
\Delta_6(\nu):=-\frac{7}{64}\sigma^6-C_6,
\qquad
\Delta_8(\nu):=C_8-\frac{23}{256}\sigma^8,
\]
and set
\[
\kappa:=\|Q_2\|_{L^2(\nu)}^2.
\]
Then
\[
\Delta_6(\nu)
=\frac{\sigma^6}{32}\beta_3^2+\frac{\sigma^6}{8}\kappa,
\]
and
\begin{align}
\Delta_8(\nu)
&=
\frac{13\sigma^2}{8}\Delta_6(\nu)
\notag\\
&\quad
+\frac{\sigma^8}{256}\left(
12\|Q_3\|_{L^2(\nu)}^2
+32\kappa^2
+\frac{61}{4}\beta_3^2\kappa
+2\beta_3^4
\right). \label{eq:defect-amplification}
\end{align}
In particular,
\[
\Delta_8(\nu)\ge \frac{13\sigma^2}{8}\Delta_6(\nu),
\]
and equality holds if and only if
\[
\nu=\frac12\delta_{\bar\gamma+\sigma}
+\frac12\delta_{\bar\gamma-\sigma}.
\]
\end{theorem}

\begin{proof}
Equation~\eqref{eq:defect-c6} gives
\[
\Delta_6(\nu)
=-\frac{\sigma^6}{64}\left(7+\frac{64}{\sigma^6}C_6\right)
=\frac{\sigma^6}{32}\beta_3^2+\frac{\sigma^6}{8}\kappa.
\]
Likewise, subtracting $23$ from \eqref{eq:defect-c8} yields
\[
\Delta_8(\nu)
=\frac{\sigma^8}{256}\Bigl(
12\|Q_3\|_{L^2(\nu)}^2
+32\kappa^2
+\frac{61}{4}\beta_3^2\kappa
+52\kappa
+2\beta_3^4
+13\beta_3^2
\Bigr).
\]
By the preceding formula for $\Delta_6(\nu)$,
\[
\frac{13\sigma^2}{8}\Delta_6(\nu)
=\frac{\sigma^8}{256}\bigl(52\kappa+13\beta_3^2\bigr),
\]
and substituting this identity proves \eqref{eq:defect-amplification}.

If equality holds in the lower bound, then every non-negative term in
\eqref{eq:defect-amplification} vanishes.  In particular,
$\beta_3=0$ and $\kappa=0$, hence $Q_2=X^2-1$ vanishes
$\nu$-almost everywhere.  Since $\EE[X]=0$ and $\EE[X^2]=1$, it follows
that $X=\pm 1$ with equal weights, which is exactly the stated
two-point law.  The converse is immediate.
\end{proof}

\begin{theorem}[Quantitative rigidity towards the symmetric two-point law]%
\label{thm:poisson-defect-stability}
Under the hypotheses of Theorem~\ref{thm:poisson-kl-eighth}, let
$\nu_c$ be the centred law of $\Delta=\gamma-\bar\gamma$ and set
\[
\rho_\sigma:=\frac12(\delta_{-\sigma}+\delta_\sigma).
\]
Then
\[
W_2(\nu_c,\rho_\sigma)
\le
4\,\sigma^{-1/2}\Delta_6(\nu)^{1/4}
+8\,\sigma^{-2}\Delta_6(\nu)^{1/2}.
\]
In particular, for fixed $\sigma$,
\[
W_2(\nu_c,\rho_\sigma)=O\!\bigl(\Delta_6(\nu)^{1/4}\bigr)
\qquad\text{as }\Delta_6(\nu)\downarrow 0.
\]
\end{theorem}

\begin{proof}
Keep the normalized variable $X=(\gamma-\bar\gamma)/\sigma$, and define
\[
Y:=\mathbf{1}_{\{X\ge 0\}}-\mathbf{1}_{\{X<0\}}\in\{-1,1\}.
\]
For every real $x$ one has
\[
|x-\operatorname{sgn}(x)|\le |x^2-1|,
\]
with the convention $\operatorname{sgn}(0)=1$.  Hence
\[
\EE[(X-Y)^2]\le \EE[(X^2-1)^2].
\]
Since $X^2-1=\beta_3X+Q_2$, the elementary bound
$(a+b)^2\le 2a^2+2b^2$ gives
\[
\EE[(X^2-1)^2]
\le 2\beta_3^2\EE[X^2]+2\|Q_2\|_{L^2(\nu)}^2
=2\beta_3^2+2\kappa.
\]
The formula for $\Delta_6(\nu)$ from
Theorem~\ref{thm:poisson-defect-amplification} therefore yields
\[
\EE[(X-Y)^2]\le \frac{64}{\sigma^6}\Delta_6(\nu).
\]
Since $\EE[X]=0$,
\[
|\EE[Y]|=|\EE[Y-X]|
\le \EE[|Y-X|]
\le \EE[(Y-X)^2]^{1/2}
\le \frac{8}{\sigma^3}\Delta_6(\nu)^{1/2}.
\]
Let $Z$ be a symmetric Rademacher random variable.  Because both
$\mathcal L(Y)$ and $\mathcal L(Z)$ are supported on $\{-1,1\}$,
\[
W_2(\mathcal L(Y),\mathcal L(Z))^2=2|\EE[Y]|.
\]
Thus
\[
W_2(\mathcal L(Y),\mathcal L(Z))
\le \frac{4}{\sigma^{3/2}}\Delta_6(\nu)^{1/4}.
\]
Also,
\[
W_2(\mathcal L(X),\mathcal L(Y))
\le \EE[(X-Y)^2]^{1/2}
\le \frac{8}{\sigma^3}\Delta_6(\nu)^{1/2}.
\]
Since $\nu_c=\mathcal L(\sigma X)$ and $\rho_\sigma=\mathcal L(\sigma Z)$,
the scaling property of $W_2$ and the triangle inequality give
\begin{align*}
W_2(\nu_c,\rho_\sigma)
&=\sigma W_2(\mathcal L(X),\mathcal L(Z)) \\
&\le \sigma W_2(\mathcal L(X),\mathcal L(Y))
+\sigma W_2(\mathcal L(Y),\mathcal L(Z)) \\
&\le
8\,\sigma^{-2}\Delta_6(\nu)^{1/2}
+4\,\sigma^{-1/2}\Delta_6(\nu)^{1/4},
\end{align*}
which is the stated estimate.
\end{proof}

\subsection{The mode Gram kernel and Hilbert completion}

Let
\[
\omega(\dd y):=\frac{\dd y}{\pi(1+y^2)}.
\]
For $\varepsilon\in\RR$ define the secant family
\begin{equation}\label{eq:mode-secant-family}
\Psi_\varepsilon(y):=\frac{R_\varepsilon(y)-1}{\varepsilon}
=\frac{2y-\varepsilon}{1+(y-\varepsilon)^2},
\qquad y\in\RR,
\end{equation}
where the formula remains valid at $\varepsilon=0$ and gives
\[
\Psi_0(y)=u_1(y)=\frac{2y}{1+y^2}.
\]
Since
\[
R_\varepsilon(y)\,\omega(\dd y)=P_1(y-\varepsilon)\,\dd y,
\]
one has
\[
\int_{\RR}R_\varepsilon(y)\,\omega(\dd y)=1
\qquad(\varepsilon\in\RR),
\]
and therefore $\Psi_\varepsilon\in L^2_0(\omega)$ for every
$\varepsilon\in\RR$.

\begin{theorem}[The mode Gram kernel]%
\label{thm:mode-gram-kernel}
For every $\varepsilon,\eta\in\RR$,
\begin{equation}\label{eq:mode-gram-kernel}
\langle \Psi_\varepsilon,\Psi_\eta\rangle_{L^2(\omega)}
=\frac{2}{4+(\varepsilon-\eta)^2}.
\end{equation}
Equivalently,
\begin{equation}\label{eq:mode-gram-RR}
\int_{\RR}R_\varepsilon(y)R_\eta(y)\,\omega(\dd y)
=\frac{\varepsilon^2+\eta^2+4}{(\varepsilon-\eta)^2+4}.
\end{equation}
\end{theorem}

\begin{proof}
For $\varepsilon\neq\eta$ define
\[
I(\varepsilon,\eta)
:=\frac{1}{\pi}\int_{\RR}
\frac{1+y^2}{(1+(y-\varepsilon)^2)(1+(y-\eta)^2)}\,\dd y.
\]
The integrand is $O(|z|^{-2})$ on the upper half-plane and has simple
poles at $z=\varepsilon+\ii$ and $z=\eta+\ii$.  A direct residue
computation gives
\[
\operatorname{Res}_{z=\varepsilon+\ii}
=\frac{\varepsilon(\varepsilon+2\ii)}
{2\ii(\varepsilon-\eta)(\varepsilon-\eta+2\ii)},
\]
\[
\operatorname{Res}_{z=\eta+\ii}
=\frac{\eta(\eta+2\ii)}
{2\ii(\eta-\varepsilon)(\eta-\varepsilon+2\ii)}.
\]
Hence, by the residue theorem,
\[
I(\varepsilon,\eta)
=2\ii\Bigl(
\operatorname{Res}_{z=\varepsilon+\ii}
+\operatorname{Res}_{z=\eta+\ii}
\Bigr)
=\frac{\varepsilon^2+\eta^2+4}{(\varepsilon-\eta)^2+4}.
\]
This proves \eqref{eq:mode-gram-RR} for $\varepsilon\neq\eta$, and the
diagonal case follows by continuity.

Since
\[
\int_{\RR}R_\varepsilon(y)\,\omega(\dd y)=1,
\]
one obtains
\[
\int_{\RR}(R_\varepsilon(y)-1)(R_\eta(y)-1)\,\omega(\dd y)
=\frac{2\varepsilon\eta}{4+(\varepsilon-\eta)^2}.
\]
For $\varepsilon\eta\neq 0$, division by $\varepsilon\eta$ yields
\eqref{eq:mode-gram-kernel}; the remaining cases follow by continuity in
$(\varepsilon,\eta)$.
\end{proof}

\begin{corollary}[Closed formulas for quadratic mode integrals]%
\label{cor:mode-quadratic-integrals}
For $|\varepsilon|$ small one has the $L^2(\omega)$ expansion
\[
\Psi_\varepsilon
=\sum_{n=1}^{\infty}u_n\,\varepsilon^{n-1}.
\]
Consequently, for every $p,q\ge 1$,
\[
\langle u_{2p},u_{2q}\rangle_{L^2(\omega)}
=(-1)^{p+q}2^{-(2p+2q-1)}
\binom{2p+2q-2}{2p-1},
\]
for every $p,q\ge 0$,
\[
\langle u_{2p+1},u_{2q+1}\rangle_{L^2(\omega)}
=(-1)^{p+q}2^{-(2p+2q+1)}
\binom{2p+2q}{2p},
\]
and
\[
\langle u_{2p},u_{2q+1}\rangle_{L^2(\omega)}=0.
\]
\end{corollary}

\begin{proof}
By Proposition~\ref{prop:cayley-mode-remainder},
\[
\Psi_\varepsilon
=\sum_{n=1}^{N}u_n\,\varepsilon^{n-1}
+\varepsilon^N v_N(\cdot,\varepsilon)
\]
with $v_N$ bounded uniformly on~$\RR^2$.  Since $\omega$ is a probability
measure, this implies convergence in $L^2(\omega)$ for $|\varepsilon|$
sufficiently small.  Hence
\[
\frac{2}{4+(\varepsilon-\eta)^2}
=\sum_{m,n\ge 1}
\langle u_m,u_n\rangle_{L^2(\omega)}\,
\varepsilon^{m-1}\eta^{n-1}.
\]
On the other hand,
\[
\frac{2}{4+(\varepsilon-\eta)^2}
=\frac12\sum_{\ell=0}^{\infty}
(-1)^\ell 4^{-\ell}(\varepsilon-\eta)^{2\ell}.
\]
Comparing coefficients gives the stated formulas.  Mixed parity terms
vanish because the kernel only contains even total degree.
\end{proof}

\begin{theorem}[The mode space and its RKHS completion]%
\label{thm:mode-space-rkhs}
Let
\[
L^2_0(\omega):=
\left\{
f\in L^2(\omega):\ \int_{\RR}f(y)\,\omega(\dd y)=0
\right\},
\]
and let $\mathcal H_K$ be the reproducing-kernel Hilbert space generated
by the kernel
\[
K(a,b):=\frac{2}{4+(a-b)^2}.
\]
If
\[
\mathcal S:=\overline{\operatorname{span}}\{\Psi_\varepsilon:\varepsilon\in\RR\}
\subset L^2_0(\omega),
\]
then
\[
\mathcal S=L^2_0(\omega).
\]
Moreover, the linear assignment
\[
\mathcal U_0(\Psi_\varepsilon):=K(\cdot,\varepsilon)
\qquad(\varepsilon\in\RR)
\]
extends uniquely to a unitary operator
\[
\mathcal U:L^2_0(\omega)\longrightarrow \mathcal H_K.
\]
\end{theorem}

\begin{proof}
We already noted that $\Psi_\varepsilon\in L^2_0(\omega)$, so
$\mathcal S\subset L^2_0(\omega)$.  Let
$f\in L^2_0(\omega)$ be orthogonal to~$\mathcal S$.  Then for every
$\varepsilon\in\RR$,
\[
0=\langle f,\Psi_\varepsilon\rangle_{L^2(\omega)}.
\]
If $\varepsilon\neq 0$, multiplying by~$\varepsilon$ gives
\[
0=\int_{\RR}f(y)\bigl(R_\varepsilon(y)-1\bigr)\,\omega(\dd y)
=\int_{\RR}f(y)R_\varepsilon(y)\,\omega(\dd y),
\]
because $\int f\,\dd\omega=0$.  Since
$R_\varepsilon(y)\,\omega(\dd y)=P_1(y-\varepsilon)\,\dd y$, this means
\[
0=\int_{\RR}f(y)P_1(\varepsilon-y)\,\dd y
\qquad(\varepsilon\in\RR).
\]
In other words, the Poisson regularization $T_1f=P_1*f$ vanishes
identically.  Now $f\in L^2(\omega)\subset\mathscr S'(\RR)$ as a tempered
distribution, so taking Fourier transforms gives
\[
e^{-|\xi|}\widehat f(\xi)=\widehat{T_1f}(\xi)=0.
\]
Since $e^{-|\xi|}$ never vanishes, $\widehat f=0$ in $\mathscr S'(\RR)$,
hence $f=0$.  Therefore $\mathcal S^\perp=\{0\}$ in $L^2_0(\omega)$, so
$\mathcal S=L^2_0(\omega)$.

Theorem~\ref{thm:mode-gram-kernel} shows that
\[
\langle \Psi_\varepsilon,\Psi_\eta\rangle_{L^2(\omega)}
=K(\varepsilon,\eta)
=\langle K(\cdot,\varepsilon),K(\cdot,\eta)\rangle_{\mathcal H_K},
\]
so $\mathcal U_0$ is an isometry on
$\operatorname{span}\{\Psi_\varepsilon:\varepsilon\in\RR\}$.  Since the
kernel sections span a dense subspace of~$\mathcal H_K$, the isometry
extends uniquely to a surjective isometry
$\mathcal U:L^2_0(\omega)\to\mathcal H_K$, that is, to a unitary
operator.
\end{proof}

\subsection{Centered reconstruction from Poisson data}

\begin{theorem}[Centered Poisson reconstruction by Laplace uniqueness]%
\label{thm:poisson-central-two-channel}
Let $\nu$ be a probability measure on~$\RR$ with finite mean~$\bar\gamma$.
Let $\nu_c$ denote the law of the centred variable
$\Delta:=\gamma-\bar\gamma$.  Define
\[
h_t(x):=\int_{\RR}P_t(x-\gamma)\dd\nu(\gamma),
\]
the central observation functions
\[
A(t):=\pi\,h_t(\bar\gamma),
\qquad
B(t):=\pi\,\partial_x h_t(x)\big|_{x=\bar\gamma}
\qquad(t>0).
\]
\[
H(t):=\int_t^\infty B(s)\dd s
\qquad(t>0),
\]
and let
\[
\phi_{\nu_c}(r):=\int_{\RR}e^{ir\Delta}\dd\nu(\gamma),
\qquad r\in\RR,
\]
be the centred characteristic function.  Then for every $t>0$ one has
\begin{equation}\label{eq:poisson-laplace-reconstruction}
A(t)+\ii H(t)
=\int_0^\infty e^{-tr}\phi_{\nu_c}(r)\dd r.
\end{equation}
Consequently, the pair $(A,B)$ determines the centred law~$\nu_c$
uniquely.  Equivalently, the triple $(\bar\gamma,A,B)$ determines the
original measure~$\nu$ uniquely.
\end{theorem}

\begin{proof}
By the elementary Laplace identities
\[
\int_0^\infty e^{-tr}\cos(r\Delta)\dd r=\frac{t}{t^2+\Delta^2},
\qquad
\int_0^\infty e^{-tr}\sin(r\Delta)\dd r=\frac{\Delta}{t^2+\Delta^2},
\]
one obtains
\[
A(t)
=\int_{\RR}\frac{t}{t^2+\Delta^2}\dd\nu(\gamma)
=\int_0^\infty e^{-tr}\Re\phi_{\nu_c}(r)\dd r.
\]
Also
\[
\partial_x h_t(x)\big|_{x=\bar\gamma}
=\frac{2}{\pi}\int_{\RR}\frac{t\,\Delta}{(\Delta^2+t^2)^2}\dd\nu(\gamma),
\]
so
\[
B(t)=2t\int_{\RR}\frac{\Delta}{(\Delta^2+t^2)^2}\dd\nu(\gamma).
\]
Since
\[
\int_t^\infty \frac{2s|\Delta|}{(\Delta^2+s^2)^2}\dd s
=\frac{|\Delta|}{\Delta^2+t^2}
\le \frac{1}{2t},
\]
Fubini's theorem applies and gives
\[
H(t)
=\int_{\RR}\frac{\Delta}{t^2+\Delta^2}\dd\nu(\gamma)
=\int_0^\infty e^{-tr}\Im\phi_{\nu_c}(r)\dd r.
\]
Combining the two formulas yields
\eqref{eq:poisson-laplace-reconstruction}.  By uniqueness of the Laplace
transform, $(A,H)$, and therefore $(A,B)$, determines
$\phi_{\nu_c}|_{[0,\infty)}$.  Since
$\phi_{\nu_c}(-r)=\overline{\phi_{\nu_c}(r)}$, the entire characteristic
function is determined, and the centred law follows from the standard
inversion theorem for characteristic functions
\cite{Widder1941,Lukacs1970,Feller1971}.  Adding the mean recovers~$\nu$.
\end{proof}

\begin{corollary}[The symmetric channel]%
\label{cor:poisson-symmetrization}
Under the hypotheses of Theorem~\ref{thm:poisson-central-two-channel}, the
function $A$ determines the symmetrization of the centred law~$\nu_c$.
\end{corollary}

\begin{proof}
The real part of $\phi_{\nu_c}$ is exactly the characteristic function of
the centred symmetrization
$\frac12(\nu_c+\check\nu_c)$, where $\check\nu_c(E):=\nu_c(-E)$.  Since
$A$ is the Laplace transform of $\Re\phi_{\nu_c}$, the claim follows from
Laplace inversion.
\end{proof}

\begin{corollary}[The symmetric case]%
\label{cor:poisson-symmetric-case}
Under the hypotheses of Theorem~\ref{thm:poisson-central-two-channel}, the
centred law~$\nu_c$ is symmetric if and only if $B\equiv 0$.  In that
case $A$ alone determines~$\nu_c$, and the pair $(\bar\gamma,A)$
determines~$\nu$.
\end{corollary}

\begin{proof}
The centred law is symmetric if and only if
$\Im\phi_{\nu_c}\equiv 0$.  By
\eqref{eq:poisson-laplace-reconstruction}, this is equivalent to
$H\equiv 0$, hence to $B\equiv 0$.
\end{proof}

\begin{theorem}[Observation channels as evaluation functionals]%
\label{thm:poisson-channel-rkhs}
For $t>0$ define the centred fluctuation profile
\[
\widetilde\delta_t(y):=
\frac{h_t(\bar\gamma+ty)}{g_t(\bar\gamma+ty)}-1,
\qquad y\in\RR.
\]
Then $\widetilde\delta_t\in L^2_0(\omega)$ and
\begin{equation}\label{eq:delta-t-secant}
\widetilde\delta_t(y)
=\frac{1}{t}\int_{\RR}\Delta\,\Psi_{\Delta/t}(y)\,\dd\nu(\gamma).
\end{equation}
Consequently,
\begin{equation}\label{eq:delta-t-rkhs}
(\mathcal U\widetilde\delta_t)(a)
=\frac{1}{t}\int_{\RR}\Delta\,K(a,\Delta/t)\,\dd\nu(\gamma),
\qquad a\in\RR.
\end{equation}
Moreover,
\begin{equation}\label{eq:delta-t-origin}
\widetilde\delta_t(0)=tA(t)-1,
\end{equation}
and
\begin{equation}\label{eq:delta-t-rkhs-origin}
(\mathcal U\widetilde\delta_t)(0)=2t\,H(2t)
=2t\int_{2t}^{\infty}B(s)\,\dd s.
\end{equation}
\end{theorem}

\begin{proof}
By definition,
\[
\widetilde\delta_t(y)
=\int_{\RR}\bigl(R_{\Delta/t}(y)-1\bigr)\,\dd\nu(\gamma)
=\frac{1}{t}\int_{\RR}\Delta\,\Psi_{\Delta/t}(y)\,\dd\nu(\gamma),
\]
which is \eqref{eq:delta-t-secant}.  Since
\[
\int_{\RR}\widetilde\delta_t(y)\,\omega(\dd y)
=\int_{\RR}(h_t(x)-g_t(x))\,\dd x=0,
\]
one has $\widetilde\delta_t\in L^2_0(\omega)$.  Applying the unitary map
$\mathcal U$ from Theorem~\ref{thm:mode-space-rkhs} to
\eqref{eq:delta-t-secant} yields \eqref{eq:delta-t-rkhs}.

At $y=0$ one has
\[
g_t(\bar\gamma)=\frac{1}{\pi t},
\]
so
\[
\widetilde\delta_t(0)
=\frac{h_t(\bar\gamma)}{g_t(\bar\gamma)}-1
=tA(t)-1,
\]
which proves \eqref{eq:delta-t-origin}.  Finally, evaluating
\eqref{eq:delta-t-rkhs} at $a=0$ and using
$K(0,b)=2/(4+b^2)$ gives
\[
(\mathcal U\widetilde\delta_t)(0)
=\frac{1}{t}\int_{\RR}\Delta\frac{2}{4+(\Delta/t)^2}\,\dd\nu(\gamma)
=2t\int_{\RR}\frac{\Delta}{4t^2+\Delta^2}\,\dd\nu(\gamma).
\]
The definition of $H$ in
Theorem~\ref{thm:poisson-central-two-channel} gives
\[
H(2t)=\int_{\RR}\frac{\Delta}{4t^2+\Delta^2}\,\dd\nu(\gamma),
\]
and \eqref{eq:delta-t-rkhs-origin} follows.
\end{proof}

\begin{remark}
The previous theorem recovers the centred Cauchy transform on the
imaginary axis:
\[
G_{\nu_c}(\ii t)
=\int_{\RR}\frac{1}{\ii t-\Delta}\dd\nu(\gamma)
=H(t)+\ii A(t)
\qquad(t>0).
\]
Thus the Laplace-transform formulation strengthens the earlier
Cauchy-transform identity while also isolating the loss of the translation
parameter.  In particular, the proof of
Theorem~\ref{thm:poisson-central-two-channel} is better viewed as a
Laplace-uniqueness plus characteristic-function-inversion argument than
as a Stieltjes inversion argument.
\end{remark}

\subsection{A Poisson image realization of the mode space}

\begin{proposition}[The mode Gram space as a Poisson image space]%
\label{thm:gram-space-poisson-image}
Let
\[
K(a,b):=\frac{2}{4+(a-b)^2}=\pi\,P_2(a-b),
\qquad a,b\in\RR,
\]
and let $\mathcal H_K$ be the reproducing-kernel Hilbert space generated
by~$K$.  Define
\[
\mathcal X
:=
\left\{
f\in L^2(\RR):
\int_{\RR}e^{2|\xi|}\abs{\widehat f(\xi)}^2\dd\xi<\infty
\right\},
\]
with norm
\[
\|f\|_{\mathcal X}^2
:=\frac{1}{2\pi^2}\int_{\RR}e^{2|\xi|}\abs{\widehat f(\xi)}^2\dd\xi.
\]
Then $\mathcal H_K=\mathcal X$ isometrically.  Equivalently,
\[
\mathcal H_K=\sqrt{\pi}\,T_1(L^2(\RR)),
\qquad
\|\sqrt{\pi}\,T_1 g\|_{\mathcal H_K}=\|g\|_{L^2(\RR)}
\quad(g\in L^2(\RR)).
\]
Every $f\in\mathcal H_K$ admits a unique holomorphic extension to the
strip
\[
\mathbb S:=\{z\in\CC:\ |\Im z|<1\},
\]
given by
\[
f(z)=\frac{1}{2\pi}\int_{\RR}\widehat f(\xi)e^{\ii z\xi}\dd\xi,
\qquad z\in\mathbb S,
\]
and satisfying the sharp pointwise estimate
\[
|f(x+\ii y)|\le \frac{1}{\sqrt{2(1-y^2)}}\,\|f\|_{\mathcal H_K}
\qquad(x\in\RR,\ |y|<1).
\]
\end{proposition}

\begin{proof}
The Fourier transform of the translation kernel $k(x):=2/(4+x^2)$ is
\[
\widehat k(\xi)=\pi e^{-2|\xi|}.
\]
Hence
\[
\widehat{K(\cdot,a)}(\xi)=\pi e^{-2|\xi|}e^{-\ii a\xi}.
\]
For $f\in\mathcal X$ one therefore has
\begin{align*}
\langle f,K(\cdot,a)\rangle_{\mathcal X}
&=
\frac{1}{2\pi^2}\int_{\RR}
e^{2|\xi|}\widehat f(\xi)\,
\overline{\widehat{K(\cdot,a)}(\xi)}\dd\xi \\
&=
\frac{1}{2\pi}\int_{\RR}\widehat f(\xi)e^{\ii a\xi}\dd\xi
=f(a)
\end{align*}
by Fourier inversion.  Thus $\mathcal X$ is a reproducing-kernel Hilbert
space with kernel~$K$, and the Moore--Aronszajn uniqueness theorem yields
$\mathcal X=\mathcal H_K$ \cite{Aronszajn1950}.

Now let $g\in L^2(\RR)$ and set $f:=\sqrt{\pi}\,T_1 g$.  Since
\[
\widehat f(\xi)=\sqrt{\pi}\,e^{-|\xi|}\widehat g(\xi),
\]
Plancherel gives
\[
\|f\|_{\mathcal X}^2
=\frac{1}{2\pi^2}\int_{\RR}e^{2|\xi|}\pi e^{-2|\xi|}
\abs{\widehat g(\xi)}^2\dd\xi
=\frac{1}{2\pi}\int_{\RR}\abs{\widehat g(\xi)}^2\dd\xi
=\|g\|_{L^2}^2.
\]
Conversely, if $f\in\mathcal X$, define
\[
\widehat g(\xi):=\pi^{-1/2}e^{|\xi|}\widehat f(\xi).
\]
Then $g\in L^2(\RR)$ and
$\widehat f=\sqrt{\pi}\,e^{-|\xi|}\widehat g$, so
$f=\sqrt{\pi}\,T_1 g$.  This proves the image-space description.

If $z=x+\ii y$ with $|y|<1$, Cauchy--Schwarz gives
\begin{align*}
|f(z)|
&\le \frac{1}{2\pi}
\left(\int_{\RR}e^{2|\xi|}\abs{\widehat f(\xi)}^2\dd\xi\right)^{1/2}
\left(\int_{\RR}e^{-2|\xi|}e^{-2y\xi}\dd\xi\right)^{1/2} \\
&=
\frac{1}{\sqrt{2(1-y^2)}}\,\|f\|_{\mathcal X},
\end{align*}
because
\[
\int_{\RR}e^{-2|\xi|}e^{-2y\xi}\dd\xi=\frac{1}{1-y^2}.
\]
Hence the integral defining $f(z)$ converges absolutely and locally
uniformly on $\mathbb S$, so the extension is holomorphic there.  The same
estimate shows uniqueness by continuity from the real axis.
\end{proof}

\begin{remark}
Proposition~\ref{thm:gram-space-poisson-image} places the mode space in a
classical Hardy/Poisson framework.  Under the conformal equivalence
between the strip $\mathbb S$ and the disk, the holomorphic extension and
pointwise estimate above are the standard $H^2$ Poisson-boundary
statements transported to a translation-invariant strip model; see
\cite{Duren1970,Rudin1987}.  Thus $\mathcal H_K$ is naturally interpreted
as a Hardy-type Poisson image space rather than merely as a Fourier
multiplier range.
\end{remark}

\subsection{Hardy splitting, lattice sampling, and cardinal reconstruction}

\begin{proposition}[Hardy splitting of the strip model]%
\label{prop:poisson-hardy-splitting}
Let
\[
\mathcal H_K^+
:=
\left\{
f\in\mathcal H_K:\ \widehat f(\xi)=0 \text{ for a.e. }\xi<0
\right\},
\]
\[
\mathcal H_K^-
:=
\left\{
f\in\mathcal H_K:\ \widehat f(\xi)=0 \text{ for a.e. }\xi>0
\right\}.
\]
Then
\[
\mathcal H_K=\mathcal H_K^+\oplus \mathcal H_K^-
\]
orthogonally.  Moreover, each $f_\pm\in\mathcal H_K^\pm$ extends
holomorphically to $\mathbb S$, and $\mathcal H_K^\pm$ is a
reproducing-kernel Hilbert space on~$\mathbb S$ with kernel
\begin{equation}\label{eq:hardy-splitting-kernels}
K_+(z,w)=\frac{1}{2(2-\ii(z-\overline w))},
\qquad
K_-(z,w)=\frac{1}{2(2+\ii(z-\overline w))}.
\end{equation}
For $a,b\in\RR$ one has
\begin{equation}\label{eq:hardy-splitting-real-kernel}
K(a,b)=K_+(a,b)+K_-(a,b).
\end{equation}
\end{proposition}

\begin{proof}
The norm formula from Proposition~\ref{thm:gram-space-poisson-image}
shows that $\mathcal H_K^+$ and $\mathcal H_K^-$ are closed subspaces of
$\mathcal H_K$, and they are orthogonal because their Fourier supports are
disjoint.  The decomposition
\[
\widehat f
=\widehat f\,\mathbf 1_{[0,\infty)}
+\widehat f\,\mathbf 1_{(-\infty,0]}
\]
therefore yields the asserted orthogonal sum.

For $w\in\mathbb S$, define $K_\pm(\cdot,w)$ by
\[
\widehat{K_+(\cdot,w)}(\xi)
:=
\pi e^{-2\xi}e^{-\ii\overline w\,\xi}\mathbf 1_{[0,\infty)}(\xi),
\]
\[
\widehat{K_-(\cdot,w)}(\xi)
:=
\pi e^{2\xi}e^{-\ii\overline w\,\xi}\mathbf 1_{(-\infty,0]}(\xi).
\]
Since $|\Im w|<1$, these functions belong to $\mathcal H_K^\pm$.  If
$f_+\in\mathcal H_K^+$, then
\begin{align*}
\langle f_+,K_+(\cdot,w)\rangle_{\mathcal H_K}
&=
\frac{1}{2\pi^2}\int_0^\infty
e^{2\xi}\widehat f_+(\xi)\,
\overline{\pi e^{-2\xi}e^{-\ii\overline w\,\xi}}\dd\xi \\
&=
\frac{1}{2\pi}\int_0^\infty \widehat f_+(\xi)e^{\ii w\xi}\dd\xi
=f_+(w).
\end{align*}
The same argument gives the reproducing identity for $\mathcal H_K^-$.  By
Fourier inversion,
\[
K_+(z,w)
=\frac{1}{2\pi}\int_0^\infty
\pi e^{-2\xi}e^{\ii(z-\overline w)\xi}\dd\xi
=\frac{1}{2(2-\ii(z-\overline w))},
\]
\[
K_-(z,w)
=\frac{1}{2\pi}\int_{-\infty}^0
\pi e^{2\xi}e^{\ii(z-\overline w)\xi}\dd\xi
=\frac{1}{2(2+\ii(z-\overline w))},
\]
which proves \eqref{eq:hardy-splitting-kernels}.  Adding the two kernels
on the real axis yields \eqref{eq:hardy-splitting-real-kernel}.
\end{proof}

\begin{theorem}[Lattice sampling for the integer secant orbit]%
\label{thm:poisson-lattice-sampling}
Let
\[
\mathcal H_K(\ZZ)
:=
\overline{\operatorname{span}}\{K(\cdot,n):n\in\ZZ\}
\subset \mathcal H_K,
\]
\[
\mathcal S_{\ZZ}
:=
\overline{\operatorname{span}}\{\Psi_n:n\in\ZZ\}
\subset L^2_0(\omega).
\]
Set
\begin{equation}\label{eq:lattice-sampling-constants}
A_{\ZZ}:=\frac{\pi}{\sinh(2\pi)},
\qquad
B_{\ZZ}:=\pi\coth(2\pi).
\end{equation}
Then the families $\{K(\cdot,n)\}_{n\in\ZZ}$ and
$\{\Psi_n\}_{n\in\ZZ}$ are Riesz sequences.  More precisely, for every
finitely supported scalar sequence $(c_n)_{n\in\ZZ}$ one has
\begin{equation}\label{eq:lattice-riesz-kernel}
A_{\ZZ}\sum_{n\in\ZZ}|c_n|^2
\le
\left\|
\sum_{n\in\ZZ}c_nK(\cdot,n)
\right\|_{\mathcal H_K}^2
\le
B_{\ZZ}\sum_{n\in\ZZ}|c_n|^2,
\end{equation}
\begin{equation}\label{eq:lattice-riesz-secant}
A_{\ZZ}\sum_{n\in\ZZ}|c_n|^2
\le
\left\|
\sum_{n\in\ZZ}c_n\Psi_n
\right\|_{L^2(\omega)}^2
\le
B_{\ZZ}\sum_{n\in\ZZ}|c_n|^2.
\end{equation}
If $F\in\mathcal H_K(\ZZ)$, then
\begin{equation}\label{eq:lattice-sampling-HK}
A_{\ZZ}\|F\|_{\mathcal H_K}^2
\le
\sum_{n\in\ZZ}|F(n)|^2
\le
B_{\ZZ}\|F\|_{\mathcal H_K}^2.
\end{equation}
Consequently, the restriction map
\[
R_{\ZZ}:\mathcal H_K(\ZZ)\to \ell^2(\ZZ),
\qquad
R_{\ZZ}F:=(F(n))_{n\in\ZZ},
\]
is an isomorphism of Hilbert spaces.  Equivalently, every
$\alpha\in\ell^2(\ZZ)$ admits a unique interpolant
$F_\alpha\in\mathcal H_K(\ZZ)$ satisfying
\[
F_\alpha(n)=\alpha_n
\qquad(n\in\ZZ).
\]
\end{theorem}

\begin{proof}
For finitely supported $c=(c_n)_{n\in\ZZ}$ set
\[
F_c:=\sum_{n\in\ZZ}c_nK(\cdot,n),
\qquad
\widehat c(\theta):=\sum_{n\in\ZZ}c_ne^{-\ii n\theta},
\qquad
\theta\in[-\pi,\pi].
\]
By Theorem~\ref{thm:mode-gram-kernel},
\[
\|F_c\|_{\mathcal H_K}^2
=\sum_{m,n\in\ZZ}c_m\overline{c_n}\,K(m,n)
=\sum_{m,n\in\ZZ}c_m\overline{c_n}\,\frac{2}{4+(m-n)^2}.
\]
Hence Parseval's identity on~$\ZZ$ gives
\begin{equation}\label{eq:lattice-symbol-form}
\|F_c\|_{\mathcal H_K}^2
=\frac{1}{2\pi}\int_{-\pi}^{\pi}\sigma(\theta)\,
|\widehat c(\theta)|^2\dd\theta,
\end{equation}
where
\[
\sigma(\theta):=\sum_{n\in\ZZ}\frac{2}{4+n^2}e^{-\ii n\theta}.
\]
Let $k(x):=2/(4+x^2)$.  Proposition~\ref{thm:gram-space-poisson-image}
gives $\widehat k(\xi)=\pi e^{-2|\xi|}$.  By the Poisson summation formula
\cite[Chapter~III]{Stein1971},
\[
\sigma(\theta)=\sum_{m\in\ZZ}\widehat k(\theta+2\pi m)
=\pi\sum_{m\in\ZZ}e^{-2|\theta+2\pi m|}.
\]
For $\theta\in[-\pi,\pi]$ this sum evaluates to
\begin{equation}\label{eq:lattice-symbol}
\sigma(\theta)
=\frac{\pi\cosh(2(\pi-|\theta|))}{\sinh(2\pi)}.
\end{equation}
Consequently,
\[
\frac{\pi}{\sinh(2\pi)}
\le
\sigma(\theta)
\le
\pi\coth(2\pi)
\qquad(\theta\in[-\pi,\pi]).
\]
Combining these bounds with \eqref{eq:lattice-symbol-form} and Parseval
for $\widehat c$ proves \eqref{eq:lattice-riesz-kernel}.  Since
$\mathcal U\Psi_n=K(\cdot,n)$ by
Theorem~\ref{thm:mode-space-rkhs}, \eqref{eq:lattice-riesz-secant}
follows immediately.

Now let $G$ be the Toeplitz Gram operator on $\ell^2(\ZZ)$ associated
with the matrix $(K(m,n))_{m,n\in\ZZ}$.  For finitely supported $c$ one
has $R_{\ZZ}F_c=Gc$, so
\[
\sum_{n\in\ZZ}|F_c(n)|^2
=\|Gc\|_{\ell^2}^2
=\frac{1}{2\pi}\int_{-\pi}^{\pi}\sigma(\theta)^2
|\widehat c(\theta)|^2\dd\theta.
\]
Since $A_{\ZZ}\le \sigma\le B_{\ZZ}$, comparison with
\eqref{eq:lattice-symbol-form} yields
\[
A_{\ZZ}\|F_c\|_{\mathcal H_K}^2
\le
\sum_{n\in\ZZ}|F_c(n)|^2
\le
B_{\ZZ}\|F_c\|_{\mathcal H_K}^2.
\]
Passing to limits in $\mathcal H_K(\ZZ)$ proves
\eqref{eq:lattice-sampling-HK}.

The symbol bounds imply that $G$ extends to a boundedly invertible
operator on $\ell^2(\ZZ)$.  Thus for every $\alpha\in\ell^2(\ZZ)$ there is
a unique $c\in\ell^2(\ZZ)$ with $Gc=\alpha$.  By
\eqref{eq:lattice-riesz-kernel}, the series $\sum_n c_nK(\cdot,n)$
converges in $\mathcal H_K(\ZZ)$ to some $F_\alpha$, and then
$R_{\ZZ}F_\alpha=\alpha$.  Uniqueness follows from the lower bound in
\eqref{eq:lattice-sampling-HK}.
\end{proof}

\begin{theorem}[Cardinal reconstruction in the lattice-generated strip RKHS]%
\label{thm:poisson-cardinal-reconstruction}
Under the hypotheses of Theorem~\ref{thm:poisson-lattice-sampling}, let
\[
L:=R_{\ZZ}^{-1}\delta_0\in\mathcal H_K(\ZZ).
\]
Then $L(n)=\delta_{0n}$ for every $n\in\ZZ$.  For every
$\alpha=(\alpha_n)_{n\in\ZZ}\in\ell^2(\ZZ)$, the series
\[
F_\alpha(z):=\sum_{n\in\ZZ}\alpha_nL(z-n)
\]
converges in $\mathcal H_K(\ZZ)$ and uniformly on compact subsets of the
strip $\mathbb S$, and it is the unique function in $\mathcal H_K(\ZZ)$
satisfying
\[
F_\alpha(n)=\alpha_n
\qquad(n\in\ZZ).
\]
If
\[
\widehat\alpha(\theta):=\sum_{n\in\ZZ}\alpha_ne^{-\ii n\theta},
\qquad \theta\in[-\pi,\pi],
\]
then
\begin{equation}\label{eq:cardinal-norm-formula}
\|F_\alpha\|_{\mathcal H_K}^2
=\frac{1}{2\pi}\int_{-\pi}^{\pi}
\frac{|\widehat\alpha(\theta)|^2}{\sigma(\theta)}\dd\theta,
\end{equation}
where $\sigma$ is the lattice symbol from \eqref{eq:lattice-symbol}.  If
$d\in\ell^2(\ZZ)$ is the unique sequence satisfying
\[
\widehat d(\theta)=\sigma(\theta)^{-1}
\qquad\text{for a.e. }\theta\in[-\pi,\pi],
\]
then
\[
L(z)=\sum_{n\in\ZZ}d_nK(z,n)
\]
in $\mathcal H_K$, and for almost every $\theta\in[-\pi,\pi]$ and every
$m\in\ZZ$ one has
\begin{equation}\label{eq:cardinal-kernel-fourier}
\widehat L(\theta+2\pi m)
=\frac{\pi e^{-2|\theta+2\pi m|}}{\sigma(\theta)}.
\end{equation}
\end{theorem}

\begin{proof}
For $n\in\ZZ$ define the translation operator
$\tau_nF(z):=F(z-n)$.  Since
\[
K(z-n,m)=K(z,m+n),
\]
integer translations preserve $\mathcal H_K(\ZZ)$.  By definition,
$R_{\ZZ}L=\delta_0$.  If $S_n$ denotes the shift operator on
$\ell^2(\ZZ)$ given by $(S_na)_m:=a_{m-n}$, then
$R_{\ZZ}\tau_n=S_nR_{\ZZ}$, so for every $n\in\ZZ$
one has
\[
R_{\ZZ}(\tau_nL)=S_n\delta_0=\delta_n.
\]
Hence $\tau_nL$ is the unique interpolant of $\delta_n$, and therefore
$L(n)=\delta_{0n}$.

If $\alpha$ has finite support, then
\[
R_{\ZZ}\left(\sum_{n\in\ZZ}\alpha_nL(\cdot-n)\right)
=\sum_{n\in\ZZ}\alpha_n\delta_n
=\alpha.
\]
By the uniqueness part of
Theorem~\ref{thm:poisson-lattice-sampling}, the corresponding interpolant
is exactly
\[
F_\alpha=\sum_{n\in\ZZ}\alpha_nL(\cdot-n).
\]
For general $\alpha\in\ell^2(\ZZ)$, let $\alpha^{(N)}$ be the truncation
to $\{-N,\dots,N\}$.  Then $\alpha^{(N)}\to\alpha$ in $\ell^2(\ZZ)$, and
the continuity of $R_{\ZZ}^{-1}$ implies
\[
F_{\alpha^{(N)}}=R_{\ZZ}^{-1}\alpha^{(N)}
\longrightarrow
R_{\ZZ}^{-1}\alpha=:F_\alpha
\]
in $\mathcal H_K(\ZZ)$.  Since each $F_{\alpha^{(N)}}$ is a finite
cardinal sum, the series for $F_\alpha$ converges in $\mathcal H_K(\ZZ)$.
Proposition~\ref{thm:gram-space-poisson-image} gives the uniform estimate
\[
|F(z)|\le \frac{1}{\sqrt{2(1-r^2)}}\,\|F\|_{\mathcal H_K}
\qquad(|\Im z|\le r<1),
\]
so $\mathcal H_K$-norm convergence implies uniform convergence on every
compact subset of~$\mathbb S$.  The interpolation identity for $F_\alpha$
follows from the continuity of $R_{\ZZ}$.

To prove \eqref{eq:cardinal-norm-formula}, first suppose that $\alpha$
has finite support and write
\[
F_\alpha=\sum_{n\in\ZZ}c_nK(\cdot,n).
\]
Then $Gc=\alpha$, where $G$ is the Toeplitz Gram operator from the proof
of Theorem~\ref{thm:poisson-lattice-sampling}.  In the discrete Fourier
representation this becomes
\[
\widehat\alpha(\theta)=\sigma(\theta)\widehat c(\theta)
\qquad\text{for a.e. }\theta\in[-\pi,\pi].
\]
Using \eqref{eq:lattice-symbol-form}, we obtain
\[
\|F_\alpha\|_{\mathcal H_K}^2
=\frac{1}{2\pi}\int_{-\pi}^{\pi}
\sigma(\theta)|\widehat c(\theta)|^2\dd\theta
=\frac{1}{2\pi}\int_{-\pi}^{\pi}
\frac{|\widehat\alpha(\theta)|^2}{\sigma(\theta)}\dd\theta.
\]
Since $1/\sigma\in L^\infty([-\pi,\pi])$, the right-hand side is
continuous in $\alpha\in\ell^2(\ZZ)$, so the formula extends to all
$\alpha$ by density.

Now let $d:=G^{-1}\delta_0$.  Then $R_{\ZZ}^{-1}\delta_0=L$ gives
\[
L=\sum_{n\in\ZZ}d_nK(\cdot,n)
\]
in $\mathcal H_K$, and the discrete Fourier transform diagonalizes $G$,
so
\[
\widehat d(\theta)=\sigma(\theta)^{-1}
\qquad\text{for a.e. }\theta\in[-\pi,\pi].
\]
Taking Fourier transforms on~$\RR$ and using
$\widehat{K(\cdot,n)}(\xi)=\pi e^{-2|\xi|}e^{-\ii n\xi}$, we find
\[
\widehat L(\xi)
=\pi e^{-2|\xi|}\sum_{n\in\ZZ}d_ne^{-\ii n\xi}.
\]
Writing $\xi=\theta+2\pi m$ with $\theta\in[-\pi,\pi]$ and $m\in\ZZ$,
the periodicity of the discrete Fourier series yields
\[
\sum_{n\in\ZZ}d_ne^{-\ii n(\theta+2\pi m)}
=\widehat d(\theta)
=\sigma(\theta)^{-1},
\]
which proves \eqref{eq:cardinal-kernel-fourier}.
\end{proof}

\begin{remark}
Theorems~\ref{thm:poisson-lattice-sampling}
and~\ref{thm:poisson-cardinal-reconstruction} place the integer secant
orbit inside the standard functional-analytic framework of principal
shift-invariant spaces and RKHS sampling.  The point here is that the
Poisson strip kernel carries a completely explicit lattice symbol and
therefore admits both stable sampling inequalities and an exact cardinal
interpolation formula.
\end{remark}
